\paragraph{Meta-1 (Clarify the core novelty and positioning).}

\begin{itemize}[leftmargin=1.2em, itemsep=1pt, topsep=2pt]
\item \textit{``The revision must clearly explain how counterfactual querying differs from existing constrained retrieval, vector-search, nearest-neighbor, filtering, and subset-selection pipelines and explain why counterfactual querying requires capabilities beyond a composition of existing retrieval components.''}
\end{itemize}
\smallskip

\noindent\textbf{Response.}
\hw{We thank the meta-reviewer for pointing out that the distinction from existing retrieval paradigms was not sufficiently clear. Our contribution is not another retrieval pipeline or counterfactual-generation algorithm; it is a new query abstraction. We have substantially revised the paper to make the new semantics---and why existing database abstractions and optimizers cannot naturally express or optimize it---explicit.

\textbf{Counterfactual querying differs from existing paradigms in what it asks and what it optimizes.}

\emph{Compared with constrained retrieval.} Constrained retrieval finds objects that satisfy user-specified predicates while optimizing a fixed retrieval objective. Counterfactual querying instead asks for feasible \emph{outcome-changing interventions}: the opposite-outcome predicate, admissibility constraints, intervention cost, and explanatory edit are part of the query semantics, not constraints layered onto an existing retrieval problem.

\emph{Compared with vector search.} Vector search retrieves objects closest under an embedding-space distance. Counterfactual querying does not optimize embedding similarity; it optimizes intervention cost $\omega$, which represents real-world effort, may be asymmetric or non-metric, and separates feasibility from effort by assigning infinite cost to infeasible interventions.

\emph{Compared with nearest-neighbor search.} Nearest-neighbor search ranks tuples by a predefined metric. Counterfactual querying searches only over feasible, opposite-outcome instances, $\omega(q)=\{x_i:y_i\neq y_q,\ \omega(x_q,x_i)<\infty\}$, so admissibility and outcome \emph{define} the search space rather than rank neighbors within it.

\emph{Compared with filtering.} Existing systems retrieve candidates and then apply filters. In counterfactual querying feasibility is not a post-processing step: it determines which instances are legal candidates in the first place, and therefore directly shapes query planning, indexing, and candidate generation.

\emph{Compared with subset-selection pipelines.} Subset selection diversifies or reranks an already-retrieved candidate set. Counterfactual querying defines the answer as the optimizer of Eq.~(2), where intervention cost and the user-specified set objective $\Omega$ are optimized \emph{jointly}, so the set objective is part of the query semantics rather than a downstream stage.

\textbf{Why counterfactual querying is a new primitive, not a composition of these components.}

One could assemble a pipeline of vector search, filtering, reranking, and subset selection, but such a pipeline treats these operators as independent stages. Counterfactual querying instead defines a single declarative query whose semantics are jointly determined by the opposite-outcome predicate, admissibility, intervention cost $\omega$, and set objective $\Omega$. Because these components are coupled, they must be visible to the query optimizer as first-class constructs so the system can reason holistically about access paths, constraint pushdown, index construction, candidate generation, and set optimization. A staged pipeline can only approximate this behavior; it cannot express or optimize the coupled semantics as one query. Concretely, feasibility reshapes the search space rather than post-filtering it (Eq.~1), $\omega$ is an asymmetric, possibly non-metric cost that breaks the geometric assumption of existing indexes (P2), and $\Omega$ is a coupled set objective for which top-$k$ selection is optimal for \emph{none} of the objectives---confirmed empirically in Table~\ref{tab:quality}. This is precisely why elevating counterfactual reasoning to a \emph{first-class data management primitive}---with its own declarative semantics, optimization opportunities, indexing mechanisms, and execution strategies---is our contribution. We revised Sections~1 and~5.1 accordingly and added Table~\ref{tab:positioning} to make these distinctions explicit, and we position DcCA against feasibility-aware counterfactual methods such as FACE~\cite{poyiadzi2020face}, which constructs feasible paths over observed points but remains per-query and model-centric and lacks a first-class query model, cost specification $\omega$, and set objective $\Omega$.}

\noindent\emph{Revision to Section~1:}
\begin{quote}
\textcolor{blue}{Counterfactual querying is not a composition of vector search, filtering, and subset selection. The opposite-outcome predicate and admissibility jointly define the feasible space rather than filter a similarity ranking; the cost $\omega$ is an asymmetric, possibly non-metric intervention effort rather than the index metric; and the answer optimizes a coupled set-level objective $\Omega$ rather than ranking tuples independently. Staging these as separate components---the only option in current systems---returns infeasible candidates, relies on post-hoc validation, and cannot jointly optimize $\Omega$; making $\omega$, the counterfactual predicate, and $\Omega$ first-class is precisely our contribution.}
\end{quote}

\noindent\emph{New positioning table (Section~1 / Section~5.1):}
\begin{table}[h]
\centering\small
\caption{Capabilities required by counterfactual querying vs.\ existing paradigms.}
\label{tab:positioning}
\resizebox{\linewidth}{!}{%
\begin{tabular}{lccccc}
\toprule
& Opp.-outcome & Non-metric & Feasibility & Set-level & Interpretable \\
& predicate & cost $\omega$ & reshapes space & objective $\Omega$ & edit $\Delta$ \\
\midrule
Constrained retrieval        & No      & No      & Partial & No      & No \\
Vector DB / ANN~\cite{pgvector,malkov2018efficient} & No & No & No & No & No \\
Nearest-neighbor search      & No      & No      & No      & No      & No \\
Subset selection~\cite{borodin2017max} & No & No & No & Yes & No \\
FACE~\cite{poyiadzi2020face} & Yes     & Partial & Partial & No      & Partial \\
\textbf{DcCA (ours)}         & Yes     & Yes     & Yes     & Yes     & Yes \\
\bottomrule
\end{tabular}}
\end{table}
